\documentclass[11pt, dina4, amsfont12]{article}
\topmargin=-1.0cm
\usepackage{graphicx}
\usepackage{epstopdf}
\usepackage{showlabe}
\usepackage{amssymb}
\usepackage{amsmath}
\usepackage{color}
\input psfig.sty
\usepackage[margin=1in]{geometry}

% ******************************* Author Macros *******************************

\newcommand{\pl}[2]{\frac{\partial#1}{\partial#2}}
\newcommand{\bu}{{\bf u} }
\newcommand{\p}{\partial}
\newcommand{\og}{\omega}
\newcommand{\Og}{\Omega}
\newcommand{\fl}[2]{\frac{#1}{#2}}
\newcommand{\dt}{\delta}
\newcommand{\nn}{\nonumber}
\newcommand{\ap}{\alpha}
\newcommand{\bt}{\beta}
\newcommand{\tht}{\theta}
\newcommand{\kp}{\kappa}
\newcommand{\Dt}{\Delta}
\renewcommand{\theequation}{\arabic{section}.\arabic{equation}}
\newcommand{\be}{\begin{equation}}
\newcommand{\ee}{\end{equation}}
\newcommand{\ba}{\begin{array}}
\newcommand{\ea}{\end{array}}
\newcommand{\bea}{\begin{eqnarray}}
\newcommand{\eea}{\end{eqnarray}}
\newcommand{\beas}{\begin{eqnarray*}}
\newcommand{\eeas}{\end{eqnarray*}}
\newtheorem{theorem}{Theorem}[section]
\newtheorem{lemma}{Lemma}[section]
\newtheorem{remark}{Remark}[section]
\newcommand{\bx}{{\bf x} }
\newcommand{\gm}{\gamma}
\newcommand{\vep}{\varepsilon}
\newcommand{\cphi}{\bar{\phi}}
%
\newcommand{\R}[1]{\textcolor{red}{#1}}

% *****************************************************************************

% =============================================================================
%                                   Title
% =============================================================================
\title{Effective Numerical for Capturing Main Trait in the Rare Mutation Limit}

\author{
Weijie Huang\thanks{{\tt ()}.}, \ 
Xinran Ruan\thanks{Corresponding author. School of Mathematical Sciences, Capital Normal University, Beijing 100048, People’s Republic of China {\tt (xinran.ruan@cnu.edu.cn)}.} 
}
\begin{document}
\date{}
\maketitle

\begin{abstract}

\end{abstract}

{\bf Key words. }  

% =============================================================================
%                                                              Section 1:  Introduction
% =============================================================================
\section{Introduction}
\label{section1}
% 引出模型并进行解释
Understanding how populations evolve under the joint effects of selection, mutation, and competition is a central question in evolutionary biology and mathematical ecology.
% Over the past decades, several mathematical formalisms have been proposed to describe such evolutionary dynamics.
At the microscopic level, stochastic individual-based models represent evolution as the outcome of random birth, death, and mutation events.
In the large population limit, these stochastic models can be rigorously connected to deterministic integro-differential or partial differential equations describing the dynamics of population densities \cite{}.
In our paper,  we consider a trait-structured population model, which has proven to be powerful tools for studying adaptation and selection.
In this setting, each individual carries a phenotypical trait $\theta\in \mathbb{R}$ that determines its fitness. 
And we use $n(t, x, \theta)$ to describe the population density at time $t$, position $x$ and with the trait $\theta$. 
 

% 回顾模型的研究并解释数值难点

% 相关算法研究



% =============================================================================
%                                                              Section 5:  Summary
% =============================================================================
\section{Summary}
\label{section6}

In this work, we 



% =============================================================================
%                                                            ACKNOWLEDGEMENTS
% =============================================================================
\bigskip
\noindent{\large\bf Acknowledgements}  
X. Ruan acknowledges support  from the National Natural Science Foundation of China (Grant number 12201436) and the R\&D Program of Beijing Municipal Education Commission from China, grant number KM202310028016.
W. Huang acknowledges support  from . 



% =============================================================================
%                                                                     REFERENCES
% =============================================================================


\end{document}




